% Shared preamble of the proof overview documents  <Folder>/<Name>_overview.tex.
% Each overview starts with
%     % Shared preamble of the proof overview documents  <Folder>/<Name>_overview.tex.
% Each overview starts with
%     % Shared preamble of the proof overview documents  <Folder>/<Name>_overview.tex.
% Each overview starts with
%     \input{../../scripts/overview_preamble.tex}
% (for ThinMatching/: \input{../scripts/overview_preamble.tex}) and is compiled from its own folder by
% scripts/build_overviews.py.
\documentclass[10pt,a4paper]{article}
\usepackage[T1]{fontenc}
\usepackage[utf8]{inputenc}
\usepackage[margin=2.2cm]{geometry}
\usepackage{amsmath,amssymb,amsthm}
\usepackage{array,booktabs,longtable}
\usepackage[dvipsnames]{xcolor}
\usepackage{enumitem}
\usepackage{tikz}
\usetikzlibrary{arrows.meta,positioning,shapes.multipart,fit,calc}
\usepackage[hidelinks]{hyperref}
\setlist{itemsep=2pt,topsep=4pt}
\renewcommand{\arraystretch}{1.2}

% Lean names, files and paper references
\newcommand{\lean}[1]{\textcolor{blue!50!black}{\texttt{\detokenize{#1}}}}
\newcommand{\file}[1]{\textcolor{OliveGreen!70!black}{\texttt{\detokenize{#1}}}}
\newcommand{\paperref}[1]{\textsf{#1}}
\newcommand{\hyp}[1]{\textcolor{BrickRed}{\textbf{#1}}}

% Boxes for the dependency diagrams
\tikzset{
  res/.style={draw, rounded corners, fill=blue!6, align=center, inner sep=4pt, font=\small},
  main/.style={res, fill=blue!18, very thick},
  cited/.style={res, fill=red!10, dashed},
  shared/.style={res, fill=gray!12},
  dep/.style={-{Stealth[length=2.2mm]}, thick},
}

\newcommand{\overviewtitle}[2]{%
  \begin{center}{\Large\bfseries #1}\\[4pt]{\small Proof overview of the Lean~4/Mathlib formalisation}\\[2pt]
  {\small #2}\end{center}}

% Standard paragraph describing the paper
\newcommand{\paperversion}{\noindent\textbf{Paper.} J.~Breitung, T.~Gervens, M.~Grohe.
  \emph{Graph Similarity Meets Optimal Transport}.\\ File \file{Graph_Similarity_Meets_Optimal_Transport.pdf} in the
  repository; statement, equation, page and line numbers refer to this PDF. \file{docs/NUMBERING.md} lists the Lean
  files and declarations of all numbered statements.\par\smallskip}

% (for ThinMatching/: \input{../scripts/overview_preamble.tex}) and is compiled from its own folder by
% scripts/build_overviews.py.
\documentclass[10pt,a4paper]{article}
\usepackage[T1]{fontenc}
\usepackage[utf8]{inputenc}
\usepackage[margin=2.2cm]{geometry}
\usepackage{amsmath,amssymb,amsthm}
\usepackage{array,booktabs,longtable}
\usepackage[dvipsnames]{xcolor}
\usepackage{enumitem}
\usepackage{tikz}
\usetikzlibrary{arrows.meta,positioning,shapes.multipart,fit,calc}
\usepackage[hidelinks]{hyperref}
\setlist{itemsep=2pt,topsep=4pt}
\renewcommand{\arraystretch}{1.2}

% Lean names, files and paper references
\newcommand{\lean}[1]{\textcolor{blue!50!black}{\texttt{\detokenize{#1}}}}
\newcommand{\file}[1]{\textcolor{OliveGreen!70!black}{\texttt{\detokenize{#1}}}}
\newcommand{\paperref}[1]{\textsf{#1}}
\newcommand{\hyp}[1]{\textcolor{BrickRed}{\textbf{#1}}}

% Boxes for the dependency diagrams
\tikzset{
  res/.style={draw, rounded corners, fill=blue!6, align=center, inner sep=4pt, font=\small},
  main/.style={res, fill=blue!18, very thick},
  cited/.style={res, fill=red!10, dashed},
  shared/.style={res, fill=gray!12},
  dep/.style={-{Stealth[length=2.2mm]}, thick},
}

\newcommand{\overviewtitle}[2]{%
  \begin{center}{\Large\bfseries #1}\\[4pt]{\small Proof overview of the Lean~4/Mathlib formalisation}\\[2pt]
  {\small #2}\end{center}}

% Standard paragraph describing the paper
\newcommand{\paperversion}{\noindent\textbf{Paper.} J.~Breitung, T.~Gervens, M.~Grohe.
  \emph{Graph Similarity Meets Optimal Transport}.\\ File \file{Graph_Similarity_Meets_Optimal_Transport.pdf} in the
  repository; statement, equation, page and line numbers refer to this PDF. \file{docs/NUMBERING.md} lists the Lean
  files and declarations of all numbered statements.\par\smallskip}

% (for ThinMatching/: \input{../scripts/overview_preamble.tex}) and is compiled from its own folder by
% scripts/build_overviews.py.
\documentclass[10pt,a4paper]{article}
\usepackage[T1]{fontenc}
\usepackage[utf8]{inputenc}
\usepackage[margin=2.2cm]{geometry}
\usepackage{amsmath,amssymb,amsthm}
\usepackage{array,booktabs,longtable}
\usepackage[dvipsnames]{xcolor}
\usepackage{enumitem}
\usepackage{tikz}
\usetikzlibrary{arrows.meta,positioning,shapes.multipart,fit,calc}
\usepackage[hidelinks]{hyperref}
\setlist{itemsep=2pt,topsep=4pt}
\renewcommand{\arraystretch}{1.2}

% Lean names, files and paper references
\newcommand{\lean}[1]{\textcolor{blue!50!black}{\texttt{\detokenize{#1}}}}
\newcommand{\file}[1]{\textcolor{OliveGreen!70!black}{\texttt{\detokenize{#1}}}}
\newcommand{\paperref}[1]{\textsf{#1}}
\newcommand{\hyp}[1]{\textcolor{BrickRed}{\textbf{#1}}}

% Boxes for the dependency diagrams
\tikzset{
  res/.style={draw, rounded corners, fill=blue!6, align=center, inner sep=4pt, font=\small},
  main/.style={res, fill=blue!18, very thick},
  cited/.style={res, fill=red!10, dashed},
  shared/.style={res, fill=gray!12},
  dep/.style={-{Stealth[length=2.2mm]}, thick},
}

\newcommand{\overviewtitle}[2]{%
  \begin{center}{\Large\bfseries #1}\\[4pt]{\small Proof overview of the Lean~4/Mathlib formalisation}\\[2pt]
  {\small #2}\end{center}}

% Standard paragraph describing the paper
\newcommand{\paperversion}{\noindent\textbf{Paper.} J.~Breitung, T.~Gervens, M.~Grohe.
  \emph{Graph Similarity Meets Optimal Transport}.\\ File \file{Graph_Similarity_Meets_Optimal_Transport.pdf} in the
  repository; statement, equation, page and line numbers refer to this PDF. \file{docs/NUMBERING.md} lists the Lean
  files and declarations of all numbered statements.\par\smallskip}

\begin{document}
\overviewtitle{Theorem 6: the upper bound $2$ for the edit distance}{Folder \file{GsimOT/Theorem6_EditUpperBound}}
\paperversion

\section{Statement}

\paragraph{Paper.} Theorem~6 (Section~3.1, p.~7, l.~245--247; proof in Appendix~C.2, pp.~24--25, l.~761--785):
\emph{Let $Q\in\mathrm{DS}(V,W)$ with a decomposition as in Eq.~(3). Then
$\sum_{i=1}^m\lambda_iC_{\mathrm{ed}}(\pi_i)\le2\,C_{\mathrm{ed}}(Q)$. In particular,
$d_{\mathrm{ed}}(G,H)\le2\,d^{\mathrm{OT}}_{\mathrm{ed}}(G,H)$ for all graphs $G,H$ of the same order.}
Eq.~(3) (p.~5, l.~193--195) is a Birkhoff decomposition $Q=\sum_{i=1}^m\lambda_iP_{\pi_i}$ with $\lambda_i>0$
and $\sum_i\lambda_i=1$. Normalised by $n^2$, the second statement is the first bound of Theorem~2
(p.~3, l.~119), $\delta_{\mathrm{ed}}\le2\,\delta^{\mathrm{OT}}_{\mathrm{ed}}$.

\paragraph{Lean.} For finite types $V,W$ and simple graphs $G$ on $V$, $H$ on $W$:
\begin{itemize}
\item \lean{GsimOT.theorem6}: for a finite index type $\iota$, weights $\lambda\colon\iota\to\mathbb R$ with
  \hyp{$\lambda_i\ge0$} and $\sum_i\lambda_i=1$, bijections $\pi_i\colon V\to W$, and $Q$ with
  $Q(v,w)=\sum_i\lambda_iP_{\pi_i}(v,w)$,
  \[ \textstyle\sum_i\lambda_i\,C_{\mathrm{ed}}(\pi_i)\le2\,C_{\mathrm{ed}}(Q). \]
\item \lean{GsimOT.theorem6_dEdit}: $d_{\mathrm{ed}}(G,H)\le2\,d^{\mathrm{OT}}_{\mathrm{ed}}(G,H)$, for all $G,H$.
\item \lean{GsimOT.theorem2_edit} (Theorem~2, first bound):
  $\delta_{\mathrm{ed}}(G,H)\le2\,\delta^{\mathrm{OT}}_{\mathrm{ed}}(G,H)$, for all $G,H$.
\end{itemize}
The Lean statements are slightly stronger: zero weights are allowed, $Q\in\mathrm{DS}(V,W)$ is not assumed but
follows (\lean{isDS_of_decomposition}), and the second and third statements need no assumption $|V|=|W|$
(otherwise there is no bijection and $d_{\mathrm{ed}}=0$). $C_{\mathrm{ed}}$, $d_{\mathrm{ed}}$,
$d^{\mathrm{OT}}_{\mathrm{ed}}$ are \lean{costEditBij}/\lean{costEdit}, \lean{dEdit}, \lean{dOTEdit} of
\file{Common/EditDistance.lean}, shared with Proposition~5 and Theorem~15.

\section{Where the proof lives}
\begin{longtable}{>{\raggedright\arraybackslash}p{5.2cm} >{\raggedright\arraybackslash}p{10.6cm}}
\toprule \textbf{File} & \textbf{Content} \\ \midrule\endhead \bottomrule\endfoot
\file{TraceFormula.lean} & the Frobenius inner product \lean{frob}, the products $A_GR$, $RA_H$, $RA_HR^\top$
  (\lean{adjMul}, \lean{mulAdj}, \lean{conjMul}), the algebraic facts used in the chain, and Eq.~(17)
  (\lean{costEdit_eq_frob}). \\
\file{Theorem6.lean} & Theorem~6 as a \texttt{calc} chain (\lean{theorem6}), the ``in particular''
  (\lean{theorem6_dEdit}) and Theorem~2 for $\delta_{\mathrm{ed}}$ (\lean{theorem2_edit}). \\
\file{Common/BirkhoffDecomposition.lean} & Eq.~(3): \lean{exists_birkhoff_decomposition} (from Mathlib's
  Birkhoff--von Neumann theorem) and \lean{isDS_of_decomposition}. \\
\file{Common/EditDistance.lean} & $C_{\mathrm{ed}}$, $d_{\mathrm{ed}}$, $d^{\mathrm{OT}}_{\mathrm{ed}}$,
  \lean{costEditBij_eq}, \lean{dEdit_le}, \lean{le_dOTEdit}. \\
\end{longtable}

\section{Proof layout}
\begin{center}
\begin{tikzpicture}
  \node[main] (t2) at (0,0) {Theorem 2 (first bound)\\\lean{theorem2_edit}};
  \node[main] (t6b) at (0,-1.9) {Theorem 6, ``in particular''\\\lean{theorem6_dEdit}};
  \node[shared] (bvn) at (6,-1.9) {Eq.~(3)\\\lean{exists_birkhoff_decomposition}};
  \node[main] (t6) at (-3,-4.1) {Theorem 6\\\lean{theorem6}};
  \node[shared] (ds) at (4,-4.1) {$Q\in\mathrm{DS}(V,W)$\\\lean{isDS_of_decomposition}};
  \node[res] (eq17) at (-6.6,-6.8) {Eq.~(17)\\\lean{costEdit_eq_frob}};
  \node[res] (first) at (-1.6,-6.8) {first step, (a)\\\lean{costEditBij_eq_frob_conjMul}\\\lean{frob_conjMul_eq}};
  \node[res] (alg) at (4.4,-6.8) {(b)--(g): Frobenius algebra\\\lean{frob_add_self_le}, \lean{frob_sub_self}\\\lean{frob_eq_sum_left}, \lean{mulAdj_eq_sum}\\\lean{frob_adjMul_permMat}};
  \node[shared] (ed) at (-4.1,-9.2) {edit distance (Common)\\\lean{costEditBij_eq}, \lean{dEdit_le}, \lean{le_dOTEdit}};
  \draw[dep] (t2) -- (t6b);
  \draw[dep] (t6b) -- (t6);
  \draw[dep] (t6b) -- (bvn);
  \draw[dep] (t6) -- (ds);
  \draw[dep] (t6) -- (eq17);
  \draw[dep] (t6) -- (first);
  \draw[dep] (t6) -- (alg);
  \draw[dep] (first) -- (ed);
  \draw[dep] (eq17) -- (ed);
\end{tikzpicture}
\end{center}
Arrows mean ``uses''. Blue: this folder; grey: \file{GsimOT/Common}. Nothing is cited; Eq.~(3) is derived from
Mathlib's \lean{exists_eq_sum_perm_of_mem_doublyStochastic}.

\section{The proof step by step}

\paragraph{Setting (l.~761--762).} The paper takes the random bijection $\pi$ with $\Pr(\pi=\pi_i)=\lambda_i$ and
$P:=P_\pi$, so $\mathbb E[P]=Q$. The formalisation does not model randomness: $\mathbb E[f(\pi)]$ is the sum
$\sum_i\lambda_if(\pi_i)$. In \lean{theorem6}, $X_i=A_GP_{\pi_i}$, $Y_i=P_{\pi_i}A_H$ and $C=QA_H$; the facts
$\mathbb E[P]=Q$ are \lean{mulAdj_eq_sum} ($QA_H=\sum_i\lambda_iP_{\pi_i}A_H$) and \lean{adjMul_eq_sum}
($A_GQ=\sum_i\lambda_iA_GP_{\pi_i}$), and $\mathbb E[c]=c$ uses $\sum_i\lambda_i=1$.

\paragraph{Eq.~(17) (l.~762--765).} $C_{\mathrm{ed}}(Q)=\|A_G\|_F^2+\|A_H\|_F^2-2\langle A_GQ,QA_H\rangle_F$ for
$Q\in\mathrm{DS}(V,W)$: \lean{costEdit_eq_frob}. As in the paper, \lean{abs_disc_eq} expands
$|a-b|=a^2+b^2-2ab$ for $a,b\in\{0,1\}$, the row sums reduce the first term to $\|A_G\|_F^2$, the column sums
the second to $\|A_H\|_F^2$, and the mixed term is regrouped into $\langle A_GQ,QA_H\rangle_F$.

\paragraph{The chain (pp.~24--25, l.~767--777).} The \texttt{calc} block of \lean{theorem6} has one step per line,
commented with the paper's line and label:
\begin{longtable}{>{\raggedright\arraybackslash}p{2.6cm} >{\raggedright\arraybackslash}p{6.2cm} >{\raggedright\arraybackslash}p{6.6cm}}
\toprule \textbf{Line, label} & \textbf{Paper} & \textbf{Lean} \\ \midrule\endhead \bottomrule\endfoot
l.~767 & $\mathbb E[C_{\mathrm{ed}}(\pi)]=\mathbb E\|A_G-PA_HP^\top\|_F^2$ & \lean{costEditBij_eq_frob_conjMul}
  (entries in $\{0,1\}$, so $|a-b|=(a-b)^2$; uses \lean{costEditBij_eq}) \\
l.~768, (a) & permutation invariance: $\|A_GP-PA_H\|_F^2$ & \lean{frob_conjMul_eq} \\
l.~769 & regrouping $(A_GP-QA_H)+(QA_H-PA_H)$ & \lean{sub_add_sub_cancel} \\
l.~770, (b) & $\|X+Y\|_F^2\le2\|X\|_F^2+2\|Y\|_F^2$ & \lean{frob_add_self_le} \\
l.~771, (c) & expanding the second square & \lean{frob_sub_self} \\
l.~772, (d) & $\mathbb E\langle QA_H,PA_H\rangle_F=\langle QA_H,QA_H\rangle_F$ & \lean{frob_eq_sum_right},
  \lean{mulAdj_eq_sum} \\
l.~773, (e) & $\langle QA_H,QA_H\rangle_F=\|QA_H\|_F^2$, collecting terms & \lean{ring} \\
l.~774, (f) & expanding the first square & \lean{frob_sub_self} \\
l.~775 & cancelling $\|QA_H\|_F^2$ & \lean{ring} \\
l.~776, (g) & $\mathbb E[P]=Q$, permutation invariance & \lean{frob_eq_sum_left}, \lean{adjMul_eq_sum},
  \lean{frob_adjMul_permMat}, \lean{frob_mulAdj_permMat} \\
l.~777, Eq.~(17) & $=2\,C_{\mathrm{ed}}(Q)$ & \lean{costEdit_eq_frob} with \lean{isDS_of_decomposition} \\
\end{longtable}
The paper explains steps (a)--(g) in l.~778--782.

\paragraph{Second assertion (p.~25, l.~783--785).} The paper takes a minimiser $Q$ of $C_{\mathrm{ed}}$ over
$\mathrm{DS}(V,W)$ and samples $\pi$ from its decomposition: $d_{\mathrm{ed}}\le\mathbb E[C_{\mathrm{ed}}(\pi)]\le
2C_{\mathrm{ed}}(Q)=2d^{\mathrm{OT}}_{\mathrm{ed}}$. As $d^{\mathrm{OT}}_{\mathrm{ed}}$ is defined as an infimum,
\lean{theorem6_dEdit} applies Theorem~6 to \emph{every} $Q\in\mathrm{DS}(V,W)$, with the decomposition
\lean{exists_birkhoff_decomposition}: $d_{\mathrm{ed}}\le\sum_\sigma\lambda_\sigma C_{\mathrm{ed}}(\pi_\sigma)\le
2C_{\mathrm{ed}}(Q)$ (\lean{dEdit_le}), and concludes with \lean{le_dOTEdit}. \lean{theorem2_edit} divides by
$n^2$.

\section{Assumptions and trusted base}
No cited result and no hypothesis. The Birkhoff--von Neumann theorem comes from Mathlib. The main theorems
depend only on \lean{propext}, \lean{Classical.choice}, \lean{Quot.sound} (checked by
\file{scripts/CheckAxioms.lean}).

\section{Notes on the formalisation}
\begin{itemize}
\item \textbf{Hypotheses.} Only $\lambda_i\ge0$ and $\sum_i\lambda_i=1$ are assumed, over any finite index type, so
  every decomposition as in Eq.~(3) is covered; $Q\in\mathrm{DS}(V,W)$ follows (\lean{isDS_of_decomposition}).
  The decomposition \lean{exists_birkhoff_decomposition} used for the second assertion sums over all
  permutations of $V$, some with weight $0$.
\item \textbf{Infimum instead of minimiser} and \textbf{no assumption $|V|=|W|$} in \lean{theorem6_dEdit},
  see above. (The infimum is attained, \lean{exists_dOTEdit_eq}, but this proof does not need it.)
\item \textbf{Not formalised}: Example~7 (p.~7, l.~250--259) and the closing remark of Appendix~C.2 (p.~25,
  l.~786--789; arbitrary matrices via thresholds).
\end{itemize}

\section{How to check this result}
\begin{itemize}
\item Build: \texttt{lake build GsimOT.Theorem6\_EditUpperBound.Theorem6} (or \texttt{bash build.sh} for the whole
  library, which also runs \file{scripts/CheckAxioms.lean}).
\item Axioms: \texttt{\#print axioms GsimOT.theorem6}, \texttt{GsimOT.theorem6\_dEdit},
  \texttt{GsimOT.theorem2\_edit}, \texttt{GsimOT.costEdit\_eq\_frob}; expected output: \lean{propext},
  \lean{Classical.choice}, \lean{Quot.sound}.
\item The table \file{Theorem6_correspondence.pdf} in this folder lists every paper item with its Lean
  declaration and source location.
\end{itemize}
\end{document}
